\section{Introduction}
% Lukas note:
% I like to use SCQA (Situation, Complication, Question, Answer) 'system' to write the introduction, see one of the two links below
% https://analytic-storytelling.com/scqa-what-is-it-how-does-it-work-and-how-can-it-help-me/
% https://corporatefinanceinstitute.com/resources/careers/how-to-job-guides/scqa/
% https://karrierebibel.de/scqa-methode/

% 1. Describe the Current Situtation
Extended Reality (XR) experiences have the potential to be deeply immersive and awe-inspiring, but achieving an optimal experience requires fine-tuning various settings, from brightness to field of view, and also haptic feedback, and other multisensory elements such as spatial audio. Currently, users often manually adjust these parameters, through traditional menu interfaces that closely resemble conventional desktop environments. 

% 2. What is the complication, challenge identified
% Spells the reason for acting now. It contains threats / opportunities and the hurdles that need to be overcome.
Unfortunately this introduces significant friction. Frequent interruptions, especially during initial setup, can break immersion, reduce excitement and potentially lower long-term adoption rates. Additionally, these conventional settings menus may carry a higher cost than just disrupting the immediate experience; they reposition the user into a known, age-old computing paradigm that is entirely disconnected from the immersive nature of XR. Instead of feeling like an extension of the virtual world, these menus trigger memories of traditional desktop environments, forcing users into a cognitive context switch that momentarily removes them from the XR space. This disconnect creates friction, making personalization feel like a chore rather than a seamless and intuitive part of the experience.

% 3. Question
% Asks the question how the hurdles of the Complication can be overcome. How can we prevent the threat or seize the opportunity? Also, what would be the benefits if the complication would be overcome?
Given these challenges, how can we personalize XR experiences in a way that minimizes manual configuration while preserving immersion? One promising approach is Reinforcement Learning from Human Feedback (RLHF), which allows the system to learn user preferences over time. However, RLHF presents its own obstacles, such as the need for human-provided labels and the challenge of balancing automation with user control. How can RLHF be effectively applied to enhance the usability of XR systems while ensuring a seamless and adaptive experience?

% 4. (short) Answer Teaser
% Provides the answer on how to overcome the hurdles. Explains how this will help deflect the threats or seize the opportunities.
% keep this short
One solution lies in obtaining implicit feedback from the user through physiological signals. Instead of relying on explicit user input, these signals can serve as a real-time indicator of a user's preferences, engagement, and immersion. By using these data as a reward signal for RLHF, the system can dynamically adjust XR settings to optimize the experience over time without requiring frequent manual interventions.

\subsection*{Combining Neuroadaptive Technology with Reinforcement Learning Agents for Personalized XR Experiences}
% 4. (long) Anwser with teaser image: 
% Provides the answer on how to overcome the hurdles. Explains how this will help deflect the threats or seize the opportunities.
% Now extend into a longer paragraph with a teaser image explaining the system in the following subsection

% earlier work on neuroadaptive technologies has shown the concept to show promise in desktop based recording environment.
% some work have done similar things in XR but, summary on ErrRP for neuroadaptive purposes

% previous works have focused on yes/no labels whereas here we try to move this one step further, from prediction violations yes/no to continuous preference ratings using a the linear transform on a binary classification scheme.

% and with a summary and a teaser to the research question/hypotheses

\input{input/1_teaser_figure}